

The Electric Motor Utilization index represents the founding device of the
physical-intensity tradition in capacity measurement. Unlike survey-based
approaches, which elicit managerial reports of operating rates relative to a
self-defined normal, the EMU methodology grounds the utilization ratio in an
engineering relationship between installed motor capacity and observed
electricity consumption. Table~\ref{tab:EMU_proxy_CU} documents the
methodology, data sources, and successive extensions of this lineage from
\citet{Foss1963} through \citet{Shaikh1992,Shaikh2016} to
\citet{Nikiforos2016}, who formalized the physical-intensity critique as a
micro-founded theoretical result and proposed the Average Workweek of Capital
as the contemporary analogue.

\begin{table}[!h]
	\centering
	\caption{Electric Motor Utilization (EMU) as a Proxy for Capacity
		Utilization: Physical Measures Lineage}
	\label{tab:EMU_proxy_CU}
	\small
	\begin{tabular}{p{3.2cm} p{10.5cm}}
		\toprule
		Feature & Description \\
		\midrule
		
		Time Span &
		1929--1954 (manufacturing); extended to 1899--1963 by
		\citet{Shaikh1992}. \\[4pt]
		
		Data \& Methodology &
		Calculates potential motor output (in kWh) from horsepower of
		installed electric motors $\times$ 8{,}760 hours/year $\times$
		0.746 kWh per hp-hour $\div$ 0.9 (efficiency adjustment). Actual
		electricity consumption is then compared to this theoretical
		maximum. Primary sources: \citet{USfpc1945};
		\citet{UScensus1966,UScensus1967,UScensus1975};
		\citet{USdc1986}. \\[4pt]
		
		Main Results &
		Substantial long-term increases in motor utilization from 1929 to
		1954: utilization rates rose from 15.9\% (1929) to 20.9\% (1954),
		indicating rising intensity in capital use even as labour hours
		declined \citep{Foss1963}. \\[4pt]
		
		Critique of FRB &
		Indirect. Suggests that survey-based measures overlook mechanical
		intensity and changes in input use; implicitly argues for a more
		objective, physically grounded metric capable of registering
		greater long-run variation than survey-elicited operating rates. \\[4pt]
		
		Extensions by \citet{Shaikh1992} &
		Adapts and extends the EMU method using Census data and
		\citet{Christensen1969} to build a continuous series from
		1899--1963, adjusting pre-1939 estimates and interpolating from
		industry-level motor data. Spliced with McGraw-Hill survey data to
		produce a utilization measure covering 1947--1985. The extended
		series diverges from FRB in amplitude, trend, and cyclical
		pattern, capturing episodes the FRB series suppresses
		\citep[Appendix~6.6]{Shaikh2016}. \\[4pt]
		
		Formalization by \citet{Nikiforos2016} &
		Re-introduces the Average Workweek of Capital (AWW) --- hours per
		week that capital equipment is operated --- as the contemporary
		physical-intensity analogue to the EMU ratio. Formalizes the
		critique of survey-based measures through the plant-manager
		thought experiment: because adding or dropping shifts changes the
		normal against which managers report, the survey denominator moves
		with the numerator. Stationarity of the FRB series is thereby a
		construction artefact, not a property of the underlying utilization
		dynamics. AWW-based estimates exhibit non-stationary behaviour
		consistent with the EMU lineage. \\[4pt]
		
		Primary Sources &
		\citet{USfpc1945}; \citet{UScensus1966}; \citet{UScensus1967};
		\citet{UScensus1975}; \citet{USdc1986}; \citet{Christensen1969};
		\citet{Foss1963}; \citet{Shaikh1992,Shaikh2016};
		\citet{Nikiforos2016}. \\
		
		\bottomrule
	\end{tabular}
	\smallskip
	
	\noindent\textit{Source: Own elaboration based on \citet{Foss1963},
		\citet{Shaikh1992,Shaikh2016}, and \citet{Nikiforos2016}.}
\end{table}
